% ============================================================================
% Reinforcement learning and agentic optimization for quantum gates and
% quantum sensing with electrons on helium
% Target: Phys. Rev. Applied / Quantum / PRX Quantum (adjust class options below)
% ============================================================================
\documentclass[aps,prapplied,reprint,superscriptaddress,longbibliography]{revtex4-2}
% \documentclass[aps,pra,reprint,superscriptaddress,longbibliography]{revtex4-2}

\usepackage{graphicx}
\usepackage{amsmath,amssymb}
\usepackage{bm}
\usepackage{hyperref}

\begin{document}

\title{Reinforcement Learning and Agentic Optimization for Quantum Gates and Quantum Sensing with
Electrons on Helium}
% Working title.

\author{Author list TBD}
\affiliation{Affiliation TBD}

\date{\today}

\begin{abstract}
Reinforcement learning (RL) already optimizes the double-well trap potential for entanglement-based
quantum sensing on the electrons-on-helium (eHe) platform, using a proximal-policy-optimization
(PPO) agent maximizing a Fisher-information-based reward~\cite{perruzza2026}. High-fidelity
two-qubit gates on the same platform, by contrast, are currently designed by grid search over a
low-dimensional pulse-shape family~\cite{leinonen2025}. We propose extending RL-based optimization
to gate design, and further explore a broader ``agentic'' workflow in which higher-level control
(hyperparameter choice, reward shaping, campaign design across the gate/sensing/readout stack) is
itself automated or human--AI collaborative. This is a proposal manuscript: we outline the RL
environment for gate optimization, a joint gate-and-sensing reward that could unify the two existing
use cases, and the open questions that need to be resolved before results can be reported.
\end{abstract}

\maketitle

\section{Introduction}

Most reinforcement-learning approaches to quantum control, including the PPO agent already deployed
on this platform~\cite{schulman2017proximal,perruzza2026}, are \emph{model-free}: the agent treats
the environment (here, Crank--Nicolson propagation under the CI Hamiltonian) as a black box, learning
a policy purely from the scalar reward returned by each episode, with no access to the fact that the
underlying dynamics are, in this case, already fully known. \emph{Model-aware} (or model-based) RL
instead folds an explicit model of the environment's dynamics directly into the learning algorithm --
either by differentiating through a known or learned generative model to obtain a direct gradient of
the reward with respect to the control parameters, or by using the model to generate cheap synthetic
rollouts that supplement real environment interactions -- and has been reported to reduce the number
of environment interactions needed by an order of magnitude relative to model-free RL on closely
analogous noisy, time-dependent gate-control problems, using an auto-differentiable ODE parametrized
by a learnable Hamiltonian ansatz~\cite{khalid2023samplemodelbased}. The same model-aware philosophy
underlies recent quantum-metrology work that folds a known measurement/parameter-estimation model
directly into the RL loop for Bayesian experimental design~\cite{belliardo2024model}, and is reviewed
more generally, outside the quantum setting, in Ref.~\cite{moerland2023modelbased}. Two features of
the present platform make model-aware RL an unusually good fit for the gate-design problem
(Sec.~\ref{sec:model-aware-rl}): first, unlike Ref.~\cite{khalid2023samplemodelbased}, which must
learn an initially uncharacterized Hamiltonian ansatz from system interactions, the eHe CI Hamiltonian
is already known essentially exactly from the soft-Coulomb/electrode-potential model of
Ref.~\cite{beysengulov2024}, so no system-identification step is needed; second, the Crank--Nicolson
propagator already used for the Ref.~\cite{leinonen2025} grid search is a deterministic numerical
model that is differentiable in principle, making it a natural candidate to differentiate through
directly rather than to treat as an opaque simulator -- which is exactly the sample-efficiency
bottleneck flagged as an open risk below (gate-fidelity evaluation is presumably expensive per
episode). A related but distinct strategy -- reinforcement learning \emph{from demonstration},
seeding the policy with a known-good solution (here, the erf-ramp family) rather than learning an
explicit environment model -- has also recently been shown to improve robustness and sample
efficiency for quantum gate control~\cite{li2025robust}, and remains the fallback option mentioned
below if a full model-aware architecture proves difficult to implement.

Model-free reinforcement learning has shown broad success in quantum control tasks — from generic
state-preparation and gate-design problems~\cite{bukov2018reinforcement,niu2019universal} to
Bayesian-optimal experimental design in quantum metrology~\cite{belliardo2024model} — and is
reviewed extensively for quantum technology applications in general~\cite{bukov2026reinforcement}.
On the electrons-on-helium platform specifically, an Actor--Critic PPO
agent~\cite{schulman2017proximal} already optimizes the trap electrode voltages to maximize a
composite reward built from singlet/triplet state fidelities and classical Fisher information,
for the purpose of entanglement-enhanced sensing~\cite{perruzza2026}. The two-qubit gate work on the
same platform~\cite{leinonen2025}, in contrast, optimizes a much smaller parameter set (ramp and
hold times of an erf-shaped voltage pulse) by grid search — a choice that was reasonable given the
low dimensionality of that particular pulse family, but leaves open whether a richer,
RL-discoverable pulse shape could improve gate fidelity, robustness, or speed beyond the current
best results ($F=0.999$ for $\sqrt{\text{iSWAP}}$, $F=0.996$ for CZ).

This project has two parts. First, the more concrete one: port the existing RL/PPO infrastructure
from the sensing use case to gate design, replacing the current grid search over
$(t_{\rm ramp}, t_{\rm hold})$ with a policy that controls the full time-dependent voltage
trajectory $V(\lambda(t))$, and see whether it finds higher-fidelity or more robust gates than the
erf-ramp family currently assumed. Second, the more exploratory one: since the same underlying
device (trap electrodes, Hartree/CI machinery, resonator readout) supports both gates and sensing,
we ask whether a single agentic framework — potentially a hierarchical or multi-objective one — can
jointly optimize across both tasks, or even autonomously propose which numerical experiments to run
next (an ``agentic AI'' layer on top of classical RL, in the sense of an agent that plans and
sequences its own investigation rather than only optimizing a fixed reward on a fixed task).

\section{Model}

\subsection{RL environment for gate design}
\label{sec:rl-env-gate-design}

Reusing the two-qubit gate model of Ref.~\cite{leinonen2025}: state $s_t$ = (present voltage vector
$V(t)$ or its parametrization $\lambda(t)$, elapsed pulse time, present gate-state overlaps
$U_{ij}(t)$); action $a_t$ = increment to the voltage trajectory (either a direct increment to
$V(t+\delta t)$, or coefficients of a smooth basis such as a truncated Fourier/spline expansion of
$\lambda(t)$, to keep the action space low-dimensional and the resulting pulses physically smooth);
reward = a combination of final gate fidelity $F$ (Sec.~IV of Ref.~\cite{leinonen2025}) and a speed
penalty, evaluated via the same Crank--Nicolson propagation already used for the grid-search
approach. Candidate architectures: PPO (consistent with the existing sensing-RL
infrastructure~\cite{schulman2017proximal}), or a demonstration-seeded approach that starts from the
known erf-ramp solution and lets RL refine it (in the spirit of reinforcement-learning-from-demonstration
approaches to quantum control) rather than learning a pulse shape from scratch.

\subsection{Model-aware RL for gate design}
\label{sec:model-aware-rl}

Rather than (or in addition to) the black-box PPO agent of Sec.~\ref{sec:rl-env-gate-design}
treating the Crank--Nicolson propagation of Ref.~\cite{leinonen2025} purely as a reward oracle, a
model-aware architecture would exploit the fact that the propagator itself is a known, deterministic
numerical model of the environment. Two concrete variants, in increasing order of implementation
effort:
\begin{itemize}
  \item \textbf{Model-based rollouts.} Use the known propagator to generate additional, cheap
  synthetic training episodes (in the spirit of classical Dyna-style model-based
  RL~\cite{moerland2023modelbased}) alongside real Crank--Nicolson evaluations, reducing the number
  of full-fidelity propagations needed per training epoch. Lowest implementation risk, since it wraps
  the existing propagator rather than reimplementing it.
  \item \textbf{Differentiate through the propagator.} Following the auto-differentiable-ODE approach
  of Ref.~\cite{khalid2023samplemodelbased} (there applied to NV-center and transmon gates with a
  partially unknown, learned Hamiltonian ansatz), reimplement the Crank--Nicolson step in an autodiff
  framework and backpropagate the gate-fidelity/speed reward directly through the full pulse
  trajectory to obtain an exact gradient with respect to the pulse parameters $\lambda(t)$, replacing
  (or initializing) the stochastic PPO policy gradient with this direct gradient. Because the eHe CI
  Hamiltonian is already known essentially exactly (Sec.~II of Ref.~\cite{beysengulov2024}), no
  Hamiltonian-learning step analogous to Ref.~\cite{khalid2023samplemodelbased}'s own is needed --
  only a differentiable re-implementation of the existing propagator.
\end{itemize}
Either variant directly targets the sample-efficiency concern raised for this project: gate-fidelity
evaluation via Crank--Nicolson propagation is presumably far more expensive per episode than the
Fisher-information reward already used for sensing~\cite{perruzza2026}, and
Ref.~\cite{khalid2023samplemodelbased} reports an order-of-magnitude reduction in the number of
environment interactions from an analogous model-aware treatment of noisy, time-dependent gate
control. \textbf{Open question:} backpropagating through many Crank--Nicolson steps is a long,
unrolled computation graph, and unrolled-ODE gradients are known to be prone to vanishing/exploding
behavior; whether this requires gradient checkpointing, a coarser effective pulse discretization, or
truncated backpropagation-through-time is unresolved and should be checked on a small toy version of
the Hamiltonian (see the sanity check in the open questions below) before committing to a full
implementation.

\subsection{A joint gate-and-sensing reward (exploratory)}

Because both tasks share the same trap/CI backbone, a natural extension is a multi-objective reward
\begin{equation}
R(a) = w_{\rm gate}\, R_{\rm gate}(a) + w_{\rm sense}\, R_{\rm sense}(a),
\end{equation}
where $R_{\rm sense}$ is the existing Fisher-information-based reward of
Ref.~\cite{perruzza2026} and $R_{\rm gate}$ is the fidelity/speed reward above, with weights
$w_{\rm gate}, w_{\rm sense}$ set by the application (a device optimized purely for sensing need not
also be a good gate, and vice versa — the interesting question is how much is given up by demanding
both). This would let a single agent characterize the trade-off surface between ``good qubit'' and
``good sensor'' trap configurations, rather than optimizing each in isolation as is currently done.

\subsection{Agentic layer (exploratory, least defined)}

Beyond parameter optimization within a fixed RL environment, an ``agentic'' framing would have a
higher-level controller: choose which reward to optimize, when to switch tasks, when a numerical
result warrants a follow-up simulation, or when to flag a result as inconsistent with the
established claims-discipline conventions of this project (see
\texttt{skills/literature-and-notation}). This is the least concrete part of the proposal and
needs scoping before any implementation is attempted — see open questions below.

\subsection{Open questions}

\begin{itemize}
  \item Does an RL-discovered pulse shape actually beat the erf-ramp family of
  Ref.~\cite{leinonen2025}, or does the existing grid search already saturate what's achievable
  for this Hamiltonian? (Plausible null result — state it as such if so.)
  \item What is the right action-space parametrization to keep training sample-efficient while still
  letting the agent discover genuinely new pulse shapes (not just re-derive the erf ramp)?
  \item Is the joint gate-and-sensing reward actually useful, or does it just produce a mediocre
  compromise with no practical use case? Needs a concrete motivating scenario (e.g.\ a device that
  must do both without reconfiguration) before investing in it.
  \item What, concretely, would distinguish an ``agentic'' layer from a slightly more elaborate
  hyperparameter search? This needs a sharper definition before it's a research claim rather than a
  buzzword — see Ref.~\cite{bukov2026reinforcement} for how the broader RL-for-quantum-technology
  literature currently scopes this kind of question.
  \item Does a model-aware/differentiable-propagator approach (Sec.~\ref{sec:model-aware-rl}) actually
  deliver, for this specific Hamiltonian, a sample-efficiency gain comparable to the order-of-magnitude
  reported in Ref.~\cite{khalid2023samplemodelbased}, or does the engineering cost of maintaining a
  differentiable re-implementation of the Crank--Nicolson propagator outweigh the savings in the
  number of episodes needed? This should be checked before choosing between the model-free (PPO),
  model-based-rollout, and full differentiable-propagator variants for the production implementation.
\end{itemize}

\section{Results}

Not yet available.

\section{Discussion and outlook}

If the gate-RL part succeeds, the natural next step is a systematic comparison against the grid
search of Ref.~\cite{leinonen2025} across the same fidelity/robustness metrics (swap error, leakage
error, sensitivity to ramp/hold-time deviations), reusing the numerical toolkit in
\texttt{ci-dvr-hartree-numerics}. The joint reward and agentic-layer parts should be scoped down to
a single, concrete, falsifiable claim before further development.

\begin{acknowledgments}
Placeholder.
\end{acknowledgments}

\bibliography{references}

\end{document}
